\section{All three tail-three labels with arbitrary tensor partners}
\label{sec:factorwise-three-tail}

\paragraph{Status and scope.}
HD-FT11--12, reserved question Q2dj, is a new derivation by
\texttt{/root/g34\_eightfolds}.  It uses the exact previously reviewed
three character polynomials and the published genus-growth input
isolated in HD-FT10.  The final natural-degree cutoff is qualitative:
the genus-growth constants are not evaluated. Complete author and
separate root whole-proof readings pass. Root also read the full
source/status record, entire checker and first successful output.
The retained review records the precise shared input and computation scope.

Use exactly the actual quotient and factorization (FT.1)--(FT.2):
there is an actual quotient onto \(P=A_n\) or \(S_n\), an actual
coefficient factorization \(V=W\otimes U\) up to scalars, and genuine
product-one generating inertia with
\[
 e=\dim V,\qquad \sum_i\operatorname{rank}(1-T_i)=2e+6.
 \tag{FT3.1}
\]
The partner \(U\), its degree, its field and every scalar phase are
arbitrary.  The projected permutations are not assigned twelve marks.
Let \(W\) have one of the labels
\[
 (n-3,3),\qquad(n-3,1,1,1),\qquad(n-3,2,1),
 \tag{FT3.2}
\]
or its transpose, or the irreducible restriction of one of these
ordinary representations to \(A_n\).  All natural degrees below are
at least \(64\), so the alternating restrictions are irreducible.

\subsection*{HD-FT11. Comparison with the unordered-triple permutation character}

For a permutation \(x\), put \(f=X_1(x)\), \(u=X_2(x)\), \(v=X_3(x)\)
and \(m=n-f\).  Write
\[
 F_2=\binom f2+u,\qquad F_3=\binom f3+fu+v,\qquad
 B=\binom n3,\quad C_2=\binom n2,\quad a_0=B-C_2.
 \tag{FT3.3}
\]
Thus \(F_j\) is the number of fixed unordered \(j\)-sets.
In the order of the three labels in (FT3.2), the exact characters
and dimensions are
\[
 \begin{array}{c|c|c}
 &\chi(x)&a\\ \hline
 R&F_3-F_2&a_0\\
 H&\chi_R+(f-1)(1-2u)&a_0+n-1\\
 M&2\chi_R+f-2(f-1)u-3v&2a_0+n.
 \end{array}
 \tag{FT3.4}
\]
These are precisely Garsia--Goupil Table~5.1, already used and
reviewed in HD-AT301, D-HK301 and HD-MT301.
\cite[Table~5.1, p.~31; Proposition~I.1 and proof]{HD:GarsiaGoupil2009}
The proof uses the displayed untwisted polynomials.  A transpose
multiplies character values by sign, so the absolute estimates persist.

\begin{lemma}[HD-FT11]
For every \(n\ge64\), each character in (FT3.4) has dimension
\(a\ge n^3/12\), and
\[
 \frac{|\chi(x)|}{a}\le
 \begin{cases}
 F_3/B,& f\ge n/2,\\
 F_3/B+48/n^{3/2},&f<n/2.
 \end{cases}
 \tag{FT3.5}
\]
The estimate is valid on all permutations and therefore all powers.
It bounds every target eigenspace by cyclic Fourier inversion.
\end{lemma}

\begin{proof}
The dimensions satisfy \(a\ge a_0=n(n-1)(n-5)/6\ge n^3/12\).
Also \(B\ge n^3/12\).
The key exact incidence identity is
\[
 3F_3-(n-2)F_2
  =-\frac{m f(f-1)}2+(3f-n+2)u+3v\le3v.
 \tag{FT3.6}
\]
If \(3f-n+2\le0\), the two terms before \(3v\) are nonpositive.
Otherwise \(u\le m/2\) makes their sum at most
\(m(-f^2+4f-n+2)/2\le0\), because
\(-f^2+4f-n+2=6-n-(f-2)^2\).
Since \(C_2/B=3/(n-2)\), (FT3.6) gives
\[
 \frac{\chi_R}{a_0}-\frac{F_3}{B}
 \le\frac{C_2v}{a_0B}\le\frac{24}{n^3}.
 \tag{FT3.7}
\]
Here \(v\le n/3\) and \(C_2\le n^2/2\).

Put
\[
 \delta_H=(f-1)(1-2u),\qquad
 \delta_M=f/2-(f-1)u-3v/2.
 \tag{FT3.8}
\]
For all \(f\), including \(f=0\), one has
\(\delta_H\le n-1\) and \(\delta_M\le n/2\).
Indeed for \(f\ge1\) discard their nonpositive terms, and for
\(f=0\) use \(2u\le n\).
Applying (FT3.7) to \(\chi_H=\chi_R+\delta_H\) and
\(\chi_M/2=\chi_R+\delta_M\) proves the common upper bound
\[
 \frac{\chi(x)}a\le\frac{F_3}{B}
                      +\frac{24}{n^3}+\frac{12}{n^2}
                  \le\frac{F_3}{B}+\frac{36}{n^2}.
 \tag{FT3.9}
\]

We also need negative character values.  Directly from the polynomials,
\[
 \chi_R=\frac{f(f-1)(f-5)}6+(f-1)u+v\ge-n/2,\qquad
 \chi_M=\frac{f(f-2)(f-4)}3-v\ge-n/3.
 \tag{FT3.10}
\]
For the first inequality, \(f=0\) gives \(-u+v\ge-n/2\);
the polynomial at integers \(f\ge1\) is at least \(-2\).
For the second, the only negative polynomial value at a nonnegative
integer is \(-1\) at \(f=3\), where \(v\le(n-3)/3\).

At \(f=0\), \(\chi_H=u+v-1\ge-1\).  For \(f\ge1\),
using \(u\le(n-f)/2\) gives
\[
 \chi_H\ge \frac{f^3}{6}-\frac{f^2}{2}-\frac{fn}{2}-1
       \ge\frac{f^3}{12}-\frac{fn}{2}-\frac{11}{3}
       \ge-\frac{\sqrt2}{3}n^{3/2}-\frac{11}{3}
       \ge-n^{3/2}.
 \tag{FT3.11}
\]
For the middle bounds,
\(f^3-6f^2+32=(f-4)^2(f+2)\ge0\), and minimizing
\(f^3/12-fn/2\) on \(f\ge0\) gives its value
\(-\sqrt2\,n^{3/2}/3\) at \(f=\sqrt{2n}\).
Equivalently that minimum follows by factoring
\(z^3-3z+2=(z-1)^2(z+2)\).
The last bound holds for \(n\ge4\).
Together, (FT3.9)--(FT3.11) and \(a\ge n^3/12\) give
the second case of (FT3.5), with the deliberately coarse constant \(48\).

It remains to prove the error-free first case.  Suppose \(f\ge n/2\).
The same exact incidence identity, now using \(u\le m/2\) and
\(v\le m/3\), gives
\[
 {\cal D}:=B F_2-C_2F_3
  \ge \frac{n(n-1)m}{12}\,K,\qquad
             K=f^2-4f+n-4\ge n^2/4-n-4.
 \tag{FT3.12}
\]
Every nonfixed three-set meets the moved letters, so
\[
 B-F_3\le m\binom{n-1}{2}.
 \tag{FT3.13}
\]
For \(n\ge64\),
\[
 nK-6(n-1)(n-2)
 \ge n^3/4-7n^2+14n-12>0.
 \tag{FT3.14}
\]
It follows that
\({\cal D}\ge(n-1)(B-F_3)\).
For the hook label,
\[
 B\chi_H-(a_0+n-1)F_3
 =-{\cal D}+B\delta_H-(n-1)F_3
 \le-{\cal D}+(n-1)(B-F_3)\le0.
 \tag{FT3.15}
\]
The same calculation works for \(\chi_M/2\), replacing \(n-1\)
by \(n/2\le n-1\), and for \(\chi_R\) it is immediate.

Finally the displayed untwisted characters are nonnegative in this
range.  For \(R\) this follows from \(f\ge32>5\).
For \(M\), \(f(f-2)(f-4)\ge f^3/4\ge n^3/32>n\).
For \(H\), the first lower bound in (FT3.11) is positive, because
\[
 f(f^2/6-f/2-n/2)-1
 \ge n^3/48-3n^2/8-1>0.
 \tag{FT3.16}
\]
Thus the error-free signed inequalities are also absolute inequalities.
This proves (FT3.5) completely.
\end{proof}

\subsection*{HD-FT12. A qualitative exclusion using the actual triple quotient}

Use precisely the published constants \(c,d>0\) and optional
unevaluated initial threshold \(n_0\ge20\) in (FT4.11).
For \(n\ge\max\{n_0,4,2d\,3^{15}/c\}\), the actual unordered-triple
quotient has genus
\[
                         g_3>cn^2/12.
 \tag{FT3.17}
\]
This is HD-FT10's direct consequence of Monderer--Neftin
Theorem~2.5, proof equation~(2.2), valid also for its exceptional
ramification types.  No numerical values of these constants, or
independent proof of the cited precursor, are asserted.

Define the finite but unevaluated integer
\[
 N^{(3)}=\left\lceil
 \max\left\{n_0,1024,\frac{2d\,3^{15}}c,
               \left(\frac{960}c\right)^2,\sqrt{\frac{168}c}\right\}
                  \right\rceil.
 \tag{FT3.18}
\]

\begin{theorem}[HD-FT12]
No actual source (FT3.1) with any label (FT3.2), its transpose,
or its ordinary alternating restriction exists for \(n>N^{(3)}\).
The conclusion permits arbitrary factorwise tensor partners, scalars,
and coefficient fields.  It is a qualitative natural-degree bound.
\end{theorem}

\begin{proof}
First, for every permutation and \(n\ge64\),
\[
 \frac{F_3}{B}\le \frac fn+\frac7n.
 \tag{FT3.19}
\]
Indeed \(F_3\le f^3/6+f(n-f)/2+n/3\).  Put \(t=f/n\).
The inequality
\((1-1/n)^{-1}(1-2/n)^{-1}\le1+6/n\), and \(B\ge n^3/12\),
give
\[
 \frac{F_3}{B}
 \le t^3+\frac6n(t^3+t-t^2)+\frac4{n^2}
 \le t^3+\frac7n\le t+\frac7n.
 \tag{FT3.20}
\]
Here \(t^3+t-t^2\le1\) for \(0\le t\le1\).

Call \(x_i\) bad when \(f(x_i)<n/2\), and put
\(\varepsilon=48/n^{3/2}\).  Every power of a nonbad image
is nonbad.  Fourier inversion, tensor spectral convolution and
(FT3.5) give
\[
 \frac{\operatorname{rank}(1-T_i)}e
 \ge\frac{\nu(W(x_i))}{a}
 \ge\frac{\operatorname{ind}(x_i,\binom{[n]}3)}B
                  -\varepsilon\,{\bf1}_{\{i\ {\rm bad}\}}.
 \tag{FT3.21}
\]
For a bad image, using (FT3.19) in the same average gives
\[
 \frac{\nu(W(x_i))}a
 \ge\frac{\operatorname{ind}(x_i)}n-\frac7n-\varepsilon
 >\frac14-\frac7n-\frac{48}{n^{3/2}}>\frac15.
 \tag{FT3.22}
\]
The last strict inequality follows at \(n=1024\) from
\(7/1024+48/32768=17/2048<1/20\), and improves with \(n\).
Since \(e\ge a\ge n^3/12\), the actual total is
\(2+6/e\le2+72/n^3<11/5\).
Consequently there are at most ten bad entries.
This count uses neither the number of projected marks nor a
separate source assertion for \(W\).

The action of \(P\) on unordered three-sets is transitive.
Its actual connected quotient curve has degree \(B\) and genus \(g_3\).
Riemann--Hurwitz therefore gives the exact identity
\[
 \sum_i\operatorname{ind}(x_i,\binom{[n]}3)=2B-2+2g_3.
 \tag{FT3.23}
\]
Sum (FT3.21) and use (FT3.17).  The threshold makes
\(cn^2/6-2>0\), and \(B\le n^3/6\).  Hence
\[
 \begin{split}
 2+\frac6e
 &>2+\frac{cn^2/6-2}{B}-\frac{480}{n^{3/2}}\\
 &\ge2+\frac cn-\frac{12}{n^3}-\frac{480}{n^{3/2}}
 \ge2+\frac{c}{2n}-\frac{12}{n^3}
 >2+\frac{72}{n^3}\ge2+\frac6e.
 \end{split}
 \tag{FT3.24}
\]
The penultimate bounds use
\(\sqrt n>960/c\) and \(n^2>168/c\).
This contradiction proves the theorem.
\end{proof}

\paragraph{What is and is not bounded.}
Together with HD-FT01--10, this leaves only ordinary tail sizes
at most two as potentially unbounded within the stated actual
quotient and factorwise tensor setup.  It does not give numerical
values for the finitely many remaining degree intervals, nor
enlarge the effective FQ bound without evaluating its additional
genus constants.  Actual quotient maps and tensor factorizations
remain necessary; this result does not infer them from an abstract
character or a proposed branch spectrum.
